\section{La qualité des stages en ingénierie}\label{la-qualite-des-stages-en-ingenierie}

\stanceinfo{\textbf{Congrès de la FCEG 2018} [2018-01-07]}{2024-01-02}

\officialstance{La Fédération canadienne étudiante de génie estime que la qualité des programmes de stages en ingénierie au Canada est inégale et souvent inférieure aux normes, et que les programmes d'ingénierie ont la responsabilité de revoir leurs pratiques afin d'offrir aux étudiants un meilleur rapport qualité-prix pour leurs frais de programme.}

\stanceheading{La position des étudiants}
\begin{itemize}
 \item La satisfaction des étudiants à l'égard de la qualité des programmes de stages est faible dans l'ensemble du Canada, ce qui traduit une insatisfaction générale à l'égard de la valeur des services de stages offerts par les écoles d'ingénieurs, notamment en ce qui a trait aux possibilités, aux ressources fournies et à l'accessibilité.
 \item L'absence de réglementation des programmes de stages en génie entraîne une plus grande variation et des normes moins élevées pour les services offerts aux étudiants, comparativement à d'autres domaines de la formation en génie.
 \item Les lacunes des programmes de stages varient considérablement d'un établissement à l'autre, ce qui signifie qu'une réponse efficace aux problèmes liés aux programmes de stages doit être apportée à l'échelle locale.
\end{itemize}

\stanceheading{Les recommandations aux partenaires, aux parties intéressées et aux autres entités}
\begin{itemize}
 \item La satisfaction des étudiants à l'égard de la qualité des programmes de stages est faible dans l'ensemble du Canada, ce qui traduit une insatisfaction générale à l'égard de la valeur des services de stages offerts par les écoles d'ingénieurs, notamment en ce qui a trait aux possibilités, aux ressources fournies et à l'accessibilité.
 \item L'absence de réglementation des programmes de stages en génie entraîne une plus grande variation et des normes moins élevées pour les services offerts aux étudiants, comparativement à d'autres domaines de la formation en génie.
 \item Les lacunes des programmes de stages varient considérablement d'un établissement à l'autre, ce qui signifie qu'une réponse efficace aux problèmes liés aux programmes de stages doit être apportée à l'échelle locale.
\end{itemize}

\stanceheading{Ce que fait la FCEG}
\begin{itemize}
 \item La FCEG a fait de la qualité des stages un domaine d'intérêt pour son Enquête nationale des étudiants de 2017 et ses recherches connexes, afin de déterminer l'ampleur des problèmes liés aux programmes de stages dans l'ensemble du Canada.
 \item La FCEG a inclus des questions sur les stages et le développement professionnel dans l'enquête nationale de 2023 afin de déterminer les exigences relatives aux stages et les ressources disponibles dans les établissements.
\end{itemize}

\stanceheading{Ce que la FCEG envisage de faire}
\begin{itemize}
 \item La FCEG travaillera avec ses partenaires pour créer des ressources afin de partager des informations de base sur les pratiques des programmes d'ingénierie à travers le Canada, et pour plaider auprès des facultés en faveur de changements efficaces.
 \item La FCEG étudiera les mérites d'un plaidoyer auprès de l'Enseignement coopératif et l'apprentissage en milieu de travail Canada pour améliorer leur agrément des programmes de stages en ingénierie.
 \item La FCEG rapportera aux DDIC les données recueillies dans le cadre de l'enquête nationale sur les stages, afin de fournir une orientation à leurs programmes individuels.
 \item La FCEG travaillera avec les organisations régionales (WESST, CAÉGO, CRÉIQ, CÉGA) afin d'obtenir de l'information au niveau régional sur les pratiques et les politiques communes en matière de stages en génie.
\end{itemize}

\newpage
